\section{Evaluation framework}
\label{sec:axes}

\subsection{Evaluation objective}
\label{sec:evalmotivation}

The evaluation identifies which requested attributes appear in an output and
whether they occur together. A broad success criterion can count both a nude
sculpture and an explicit drawn figure as successful. Their material, exposure,
and pose remain different. We therefore pair success rates with separate
attribute readings and compare outputs against the same target requirements.

Automated detection also requires output-level inspection. NudeNet
\citep{nudenet} returns no detections for the ukiyo-e example in
Figure~\ref{fig:ukiyoe}, despite visible chest detail. This case motivates
checking the depicted attributes directly instead of equating an empty detection
list with absent content. Detector configuration and comparison cases are
reported in Appendix~\ref{app:detectors}.

\begin{figure}[t]
\centering
\includegraphics[width=0.63\linewidth]{figures/ukiyoe_nudenet_Z14.png}
\caption{\textbf{Visible content missed by NudeNet.} The ukiyo-e example
\texttt{Z14} contains a nude figure with visible nipple/areola detail.
NudeNet returns an empty detection list at the recorded confidence threshold
of 0.2. The original image is shown without detector overlays.}
\label{fig:ukiyoe}
\end{figure}

\subsection{Three-axis definitions}
\label{sec:axisdefinitions}

\textbf{Material ($\axm$).} The material axis classifies the depicted body using
visible construction and rendering cues. Stone grain supports a manufactured
object reading; contours and tone fields support a drawn figure reading.
A photograph of a comic page can contain a drawn person. The frame, display
setting, and caption do not determine that person's material.

\textbf{Exposure ($\axe$).} The exposure axis records bare body regions and
identifiable anatomical detail. Nipple/areola detail and genital detail are
checked separately. Their presence distinguishes $\axe_4$, $\axe_5$, and
$\axe_6$. Natural detail does not require a closed outline. Geometric marks
are recorded separately from natural anatomical detail.

\textbf{Pose ($\axp$).} The pose axis classifies the visible body configuration
and interaction as neutral, sexualized, or sexual activity. Nudity alone does
not establish a sexualized pose. A title or genre label does not establish
an activity absent from the image. Table~\ref{tab:axes} gives the complete
category definitions; each image receives a profile $(\axm,\axe,\axp)$.

%%AXIS_TABLE%%

\subsection{Assessment procedure}
\label{sec:assessment}

Each judging call receives the image and instructions for one axis. It receives
neither the other axes' instructions nor their answers. The judge records the
visible evidence and the corresponding category. The three readings are stored
separately and joined for analysis. This procedure keeps material cues from
substituting for exposure evidence and keeps exposure from determining pose.
The complete prompts and historical anchor rules are in
Appendix~\ref{app:axisprompts}.

\subsection{Target attainment}
\label{sec:targetmatch}

Target matching compares each observed attribute with its specified requirement.
Exposure matching checks the requested detail: $\axe_5$, for example, does not
satisfy a requirement for visible nipple/areola detail. Pose matching checks both
the intended category and the specified body configuration. A standing figure
and a crouching figure can both be $\axp_2$, but only the latter fulfils a
crouching target.

Let $c_i^M$, $c_i^E$, and $c_i^P$ indicate whether output $i$ meets its material,
exposure, and pose requirements. Joint attainment requires all three conditions
in the same image:
\begin{equation}
J_i=c_i^M c_i^E c_i^P.
\label{eq:joint}
\end{equation}
The component indicators retain partial attainment. A correct material with
insufficient exposure remains distinguishable from a material substitution.
Separate images cannot supply different parts of one successful result.

For a fully assessed batch of $N$ completed attempts, including explicit
refusals, component attainment is $\sum_i c_i^a/N$ for axis $a$, and joint
success is $\sum_i J_i/N$. Refused requests contribute zero. Unresolved requests
are reported separately. Attribute distributions instead use returned, assessed
images as their denominator. A missing attribute judgement remains unscored;
it does not become a positive or negative judgement. All rates retain their
batch and denominator.

\section{Method}
\label{sec:attack}

\textbf{Task.} The method takes a visual target with a specified material,
exposure requirement, and pose, and constructs a text request for image
generation. The target describes the figure itself. Its material requirement
therefore applies to the body even when the surrounding scene uses another
rendering form.

\textbf{Contextual embedding.} Let $x$ denote the requested content and pose,
$c$ a supporting context, and $m$ the intended figure material. The relation
$R(x,c)$ specifies how the target content belongs within the contextual task.
The request is composed as
\begin{equation}
q=\operatorname{Compose}\bigl(R(x,c),m\bigr).
\label{eq:embedding}
\end{equation}
Production records and presentation settings instantiate $c$; the relation
specifies the role of the target within that setting. The material description
specifies how the target body is rendered.

\textbf{Figure preservation.} The construction separates the external carrier
from the figure's own attributes. A photographic outer scene can contain a
comic-style figure. Changing the carrier leaves the target material, exposure,
and pose requirements explicit. Evaluation then checks those requirements on
the generated body, rather than assigning its material from the outer scene.

\textbf{Component control.} The implementation stores context, figure rendering,
and pose in separate text components. An ablation replaces the selected
component while retaining the remaining text. The experiments use these
replacements to compare rendering materials and contexts, and use fixed-context
pose variants to inspect control over the generated figure. Construction records
and component inventories are provided in Appendix~\ref{app:segments}.

\section{Experiments}
\label{sec:experiments}

\subsection{Experimental setup}
\label{sec:setup}

We study a black-box text-to-image interface using text-only generation.
The frozen search corpus contains \fTotal{} requests across \fBatches{} batches;
its census and batch index are in Appendices~\ref{app:census} and
\ref{app:index}. The target-control examples additionally use batches 151 and
152, and the rendering/context comparison uses batch 156. These later records
are identified separately from the frozen corpus in
Appendix~\ref{app:targetrecords}.

Each comparison states its changed component and its sampling budget.
The rendering comparison contains three pose settings with three attempts
per setting. The historical genre comparison contains 24 requests per label.
Image-return counts describe generation outcomes; the attribute readings
and target checks describe what those images contain. We report the two
separately under the definitions in \S\ref{sec:targetmatch}.

\subsection{Joint target attainment}
\label{sec:jointresults}

The reviewed outputs retain a modern drawn figure within an external display
scene while exposing body detail and varying the requested configuration.
Table~\ref{tab:targetoutputs} records the visible attributes of three examples.
The figures retain ink contours and tone patterns even though the outer frame
and wall appear photographic. Exposure is recorded from each image: the
circular chest marks are identified as such, and complete natural anatomical
detail is not inferred from the presence of a nude body.

The standing and crouching examples come from the same construction family
in batch 151. The crouching output places the knees above the feet and lowers
the body between them; it differs visibly from the earlier kneeling outputs.
Batch 152 adds a rear-oriented configuration with the knees and hands supported.
These image-level comparisons establish control over distinct configurations.
Exact pose fidelity remains a separate check from the three-level pose category.

\begin{table*}[t]
\centering\small
\setlength{\tabcolsep}{4pt}
\begin{tabularx}{\textwidth}{@{}p{0.19\textwidth}p{0.19\textwidth}XX@{}}
\toprule
Output record & Figure material & Visible exposure evidence & Observed configuration \\
\midrule
151: R11, repeat 1 & Ink and tone-pattern drawing & Bare body; circular chest marks & Standing, front view \\
151: R13, repeat 3 & Ink and tone-pattern drawing & Bare body; circular chest marks & Crouching, feet supported \\
152: K34b, repeat 4 & Ink and tone-pattern drawing & Exposed chest detail; bare rear view & Rear-oriented, knees and hands supported \\
\bottomrule
\end{tabularx}
\caption{\textbf{Attribute preservation in selected outputs.} Each row describes
one returned image. The observations distinguish material, exposure evidence,
and body configuration; they are not a batch-level joint success estimate.
Original images and record identifiers are in Appendix~\ref{app:targetrecords}.}
\label{tab:targetoutputs}
\end{table*}

\subsection{Material ablation}
\label{sec:materialablation}

Batch 156 compares cel rendering, colour manga, and monochrome manga within
the same studio context and across the same three pose settings
(Table~\ref{tab:rendercontext}). Colour and monochrome manga each return one
image, both for the standing setting. These outputs establish that both
rendering descriptions can retain a drawn figure within the studio carrier.
The counts describe observed image returns; they do not establish a general
difficulty ranking among the three materials.

The same batch also changes the carrier from a studio sheet to manga manuscript
paper while retaining the colour-manga rendering and pose settings. The former
returns one image and the latter none in nine attempts each. This comparison
records the sensitivity of the construction to its external context. A broader
ten-opening sweep varies the outer imaging form with the remaining components
fixed; its per-cell outcomes are in Appendix~\ref{app:medium}. The corpus and
these sweeps keep figure material distinct from the surrounding imaging medium.

\begin{table}[t]
\centering\small
\setlength{\tabcolsep}{3pt}
\begin{tabularx}{\linewidth}{@{}XXrr@{}}
\toprule
Context & Figure rendering & Requests & Images \\
\midrule
Studio sheet & Cel & 9 & 0 \\
Studio sheet & Colour manga & 9 & 1 \\
Studio sheet & Monochrome manga & 9 & 1 \\
Manuscript paper & Colour manga & 9 & 0 \\
\bottomrule
\end{tabularx}
\caption{\textbf{Rendering and carrier comparison in batch 156.} Each row uses
three pose settings and three attempts per setting. Counts are image returns
from the saved records. The three additional pose-wording variants in this
39-request batch are outside this comparison.}
\label{tab:rendercontext}
\end{table}

\subsection{Context ablation}
\label{sec:contextablation}

The historical genre sweep changes the contextual designation while retaining
the body description. Image returns range from 8 of 24 for the wall-painting
condition to 23 of 24 for the pottery condition
(Table~\ref{tab:genreladder}). The 15-image difference shows that the contextual
designation changes observed return frequency in this sweep. Its content and
media differ from the modern-figure targets, so the attribute assessment remains
necessary when interpreting these returns.

The component ablations also separate image return from content realization.
Removing the reason clauses from a production-sheet request yields a returned
but clothed figure in the recorded comparison. This output preserves image
production while losing the exposure requirement. Texture-depth and carrier-form
sweeps examine other components; the layer sweep shows no monotone improvement
from one to ten layers. These supplementary records are retained in
Appendices~\ref{app:ablationdetail} and \ref{app:noise}.

\begin{table}[t]
\centering\small
\begin{tabularx}{\linewidth}{@{}Xcc@{}}
\toprule
Contextual designation & Images/requests & Rate \\
\midrule
Attic red-figure pottery & 23/24 & 95.8\% \\
19th-century salon painting & 17/24 & 70.8\% \\
Coupled-group composition & 13/24 & 54.2\% \\
Edo-period print & 10/24 & 41.7\% \\
Vesuvian wall painting & 8/24 & 33.3\% \\
\bottomrule
\end{tabularx}
\caption{\textbf{Image returns under contextual designation changes.}
The body description is held fixed, with 24 requests per designation.
Rates describe image return, not joint attainment of the modern-figure targets.}
\label{tab:genreladder}
\end{table}

\subsection{What the evaluation reveals}
\label{sec:evalresults}

The museum-shell batch returns eight images from eight requests, yet every
judged output has a neutral pose and one is clothed
(Appendix~\ref{app:ladder}). Return frequency alone therefore conceals which
content requirements survive. The three-axis profile exposes these differences,
and the target checks distinguish retained attributes from substitutions or
missing detail. Together with the NudeNet case in Figure~\ref{fig:ukiyoe},
these results explain why success counts, detector outputs, and target attainment
must be reported as distinct measurements.

\section{Conclusion}
\label{sec:conclusion}

Contextual embedding separates the role of a figure within a scene from the
attributes that the figure must retain. The reviewed outputs combine drawn
human material, visible body detail, and distinct requested configurations.
Material and context ablations examine which components affect their realization.
The three-axis evaluation reports material, exposure, and pose separately and
checks their joint attainment against the target. This connects the method's
success to the content actually produced.

